\chapter{Discussion and Threats to Validity}\label{discussion-and-threats-to-validity}

\section{Interpretation and Practical Implications}\label{interpretation-and-practical-implications}

Adding pipeline stages does not consistently improve verification. Raw
requirements with all screenshots produce the highest accuracy in the primary
Flash-Lite matrix. Lexical top-4 omits decisive information and saves almost no
money under the implemented batching strategy. Automatic decomposition
improves macro-F1 under restricted evidence, but all-screenshot accuracy remains
unchanged and false fulfillment increases. Requirement understanding and
screenshot selection consequently need separate evaluation.

The evidence-first representation primarily improves inspectability. It records
selected screenshots, supported and unresolved claims, uncertainty reasons, and
localized evidence. A reviewer can then separate retrieval errors from visual
interpretation and label-policy disagreements. Direct whole-flow prompting and
staged verification can therefore differ in label quality and diagnostic
detail.

For a practical deployment, the results favor an adaptive rather than fixed
retrieval policy. Complete flows are appropriate while flows remain short
enough for the model context and budget. Retrieval becomes useful when flows
grow, but it should monitor evidence sufficiency and fall back to additional
screens when outcome states, transitions, or multiple obligations are
unresolved. Late-screen priors and action/result pairing are concrete
improvements suggested by the observed failures.

The evidential threshold for fulfillment is context-dependent. A
safety-critical setting may prefer more abstention, while another project may
classify missing required evidence as failure. In this study, the four label
definitions are fixed across all applications and flows, supplied to the model,
enforced by the aggregator, and evaluated with safety-oriented metrics. Other
risk policies can be compared by changing these explicit rules instead of
folding them into a binary score.

\section{Internal Validity}\label{internal-validity}

All current Mind2Web verification items were reviewed by the primary author
independently of the evaluated model predictions. The documented reference
standard includes requirement labels, claim boundaries, evidence sets, and
UI-evaluability judgments. Reported model metrics quantify agreement with this
reviewed reference standard. A separate disagreement-focused author audit
provides qualitative analysis of recurring boundary interpretations.

The requirements are derived from the same flows against which they are
evaluated. Prediction-independent manual review keeps evaluated outputs
separate from reference decisions. Separate reporting for original and
contrastive requirements, flow-level resampling, and the exploratory PURE
document comparison make the construction and evidence sources visible in the
analysis.

Model outputs are sensitive to model version, prompt text, image preparation, retries, and aggregation. Historical repository reports combine different current-per-flow runs and are useful for diagnosis but not controlled comparison. Final experiments require an immutable run manifest and exact model identifiers.

\section{Construct Validity}\label{construct-validity}

The four labels operationalize a conservative interpretation of screenshot-based verification. Other projects might define missing evidence as failure or treat visible success messages as sufficient proof of a backend outcome. The adopted label policy is therefore a deliberate construct tied to the intended safety goal rather than the only possible definition.

The label definitions specify the meaning of the four verdicts, but they do not
fully specify when visible evidence is sufficient for assigning them. When a
control, message, or interface state suggests the required behavior, the
verifier must decide whether that cue establishes the behavior or merely makes
it plausible. With only generic evidence guidance, part of this decision is
supplied by the model's implicit assumptions. The observed cross-model
disagreements are consistent with such model-dependent interpretation, even
though repeated executions of the same Gemini configuration produced stable
labels. The study retained one generic evidence policy because the 13-flow
benchmark was too narrow to derive detailed sufficiency criteria without
risking calibration to its particular applications and items.

Screenshot-step overlap is an incomplete measure of evidence quality. Human annotations may contain several valid screens, and a prediction may cite a semantically valid alternative that is absent from the reference set. Conversely, retrieving a gold step does not prove that the model used the correct region or interpretation. Manual evidence inspection should complement the automated metrics.

Region-level review introduces a related construct question. A visually tight
box is not necessarily the most useful explanation, and some claims require a
whole-screen state, multiple regions, or a transition. The grounding evaluation
therefore combines geometric measures with human relevance and sufficiency
judgments instead of defining correctness only through intersection over union.

UI evaluability is also a constructed boundary. The difference between a
routine hidden dependency behind a visible success proxy and a nontrivial
hidden property can depend on the application risk model. This boundary
therefore requires the same explicit evidence guidance as the fulfillment
labels.

Claim matching introduces additional uncertainty. Two decompositions can be semantically equivalent while using different granularity. Low claim-status performance may reflect a poor decomposition, poor text matching, or poor status prediction. The final report should separate claim-match recall from status quality on matched claims.

\section{External Validity}\label{external-validity}

Thirteen web flows are a small sample of the variety of real interfaces and requirements. The selected tasks do not establish generalization to native mobile applications, desktop software, accessibility requirements, or industrial specifications. Mind2Web trajectories reflect one recorded interaction path and may omit alternative states relevant to a requirement.

The primary benchmark also reverses the usual requirements-first development
order. Its source requirements were reconstructed from already observed web
flows, so the contract and its evidence are not independent. This construction
supports controlled evaluation of evidence interpretation, but it favors
requirements that align with established interface behavior and does not
reproduce a setting in which an independently authored requirement precedes a
design or implementation that may violate it. The study neither generated UI
drafts nor seeded implementation defects. Contrastive items introduce
unsupported and contradictory cases, but they are constructed deviations
rather than observed implementation defects. The results therefore do not
estimate defect-detection performance in requirements-driven development.

PURE broadens the source domain with realistic requirements documents. Its
results are framed as document-to-UI consistency and reported separately from
implementation-conformance results over recorded execution flows.

\section{Reliability and Reproducibility}\label{reliability-and-reproducibility}

Model APIs can return malformed responses, change behavior between versions, or fail transiently. Older flow runs contained API fallbacks, and cached responses may hide differences between repeated executions. Final reporting should include retries, failures, fallbacks, token counts, image counts, runtime, cache policy, and pricing assumptions.

The repository contains stale summary metrics generated against older benchmark snapshots. For example, one stored batched-top-k metric reports 31.8\% accuracy because it evaluates 201 predictions against all 258 current gold items and counts the 57 absent predictions as abstentions. Recomputed metrics restricted to the actual 201-item run give 75.1\% accuracy. This discrepancy demonstrates why each thesis table must be generated from a frozen manifest with matching prediction and gold sets.

\section{Current Scope Limitations}\label{current-scope-limitations}

The pipeline provides mature step-level evidence and implemented region
grounding whose final human evaluation is reported separately once complete.
Bounding boxes are not treated as an accuracy intervention. Screenshots cannot
establish hidden backend truth, global absence, long-term persistence, external
delivery, or complete result correctness. The completed 2x2 matrix isolates
claim and screenshot policy for one model, but it does not show that
evidence-first design uniformly improves false fulfillment. The reported
conclusions remain within these limits; evidence grounding alone does not
establish safer labels.

\section{Dataset Licensing and Artifact Availability}\label{dataset-licensing-and-artifact-availability}

All 13 primary flows belong to the Mind2Web \texttt{test\_task} split. Mind2Web is
identified by its maintainers as CC BY 4.0, but the official repository also
asks users not to redistribute unzipped test files online and not to place
benchmark data in training corpora. The public replication artifact is
therefore restricted to code, configurations, citations, and aggregate results.
It excludes original screenshots, HTML, MHTML, HAR files, traces, videos,
processed trajectories, test records, and per-item raw model interactions.
Exact test identifiers and detailed annotations remain available only for
examination under controlled access unless the maintainers approve broader
release.

PURE has a different rights limitation. Its curators collected requirements
documents from third-party Web sources and explicitly state that they are not
aware of license agreements or intellectual-property rights governing the
source requirements. Consequently, the public artifact excludes PURE PDFs, XML
files, extracted figures, substantial text passages, and per-item outputs that
reproduce those passages. It may contain document identifiers, hashes,
author-created non-textual labels, aggregate metrics, and code that requires
users to obtain PURE independently. These conservative boundaries support
replication without asserting rights over third-party source content.
